% A suitable documentclass line is the following. Change this to [final]
% for the submission version -- [draft] additionally typesets \shout{}
% comments and other draft-time markup.
\documentclass[draft]{exam-n}

\exambanner{Examination for the degree of Bachelor of Science}
\degreedescriptions{Physics 3\\Physics with Astrophysics 3}
\schoolcoursecode{PHYS3021}
\universitycoursecode{PHYS3021}
\coursetitle{Classical Mechanics}

\examdate{Monday, 8 December 2026}
\examtime{9:30 -- 11:30}

\rubric{Candidates should answer \textbf{all} questions.
  Each question carries the marks shown.}

\numquestions{3}

\usepackage{siunitx}

% exam-n.cls always renders a "no logo available" placeholder inside the
% title-block logo box when no institution logo is supplied (see \UniLogo
% in exam-n.cls) -- that text is baked into the class, not something the
% document controls via a macro. It read as a rendering bug rather than an
% intentional placeholder, so we redefine it to a blank ruled box: the
% border still signals "your logo goes here" without printing text that
% looks like a broken template.
\renewcommand{\UniLogo}[1]{\vbox to 15.5mm{%
    \hrule
    \vss\hbox to #1{\hss\hss}%
    \vss
    \hrule}}

\begin{document}
\maketitle

\begin{question}{20} \author{J. Halliday}
\part A particle of mass $m$ moves in one dimension under a restoring
force $F(x) = -kx$. Write down its equation of motion and give the
general solution for $x(t)$.
\partmarks{6}
\begin{solution}
$m\ddot{x} = -kx$, so $x(t) = A\cos(\omega t) + B\sin(\omega t)$ with
$\omega = \sqrt{k/m}$.
\end{solution}

\part The particle starts at rest at $x_0 = \SI{0.2}{\metre}$ with
$k/m = \SI{4}{\second^{-2}}$. Find $x(t)$ explicitly.
\partmarks{8}
\begin{solution}
$\omega = 2\,\si{\second^{-1}}$; with $x(0)=x_0$ and $\dot{x}(0)=0$ this
gives $x(t) = \SI{0.2}{\metre}\cos(2t)$.
\end{solution}

\part State the period of the oscillation found above.
\partmarks{6}
\begin{solution}
$T = 2\pi/\omega = \pi\,\si{\second} \approx \SI{3.14}{\second}$.
\end{solution}
\end{question}

\begin{question}{30} \author{J. Halliday}
A satellite of mass $m$ orbits a planet of mass $M$ in a circular orbit of
radius $r$.

\part Derive an expression for the orbital speed $v$ in terms of $G$, $M$
and $r$.
\partmarks{10}
\begin{solution}
Setting gravitational force equal to centripetal force,
$GMm/r^2 = mv^2/r$, so $v = \sqrt{GM/r}$.
\end{solution}

\part Derive Kepler's third law, $T^2 \propto r^3$, from the result above.
\partmarks{10}
\begin{solution}
$T = 2\pi r / v = 2\pi r^{3/2}/\sqrt{GM}$, so $T^2 = (4\pi^2/GM)\,r^3$,
i.e.\ $T^2 \propto r^3$.
\end{solution}

\part A second satellite orbits at $4r$. Find the ratio of its period to
that of the first satellite.
\partmarks{10}
\begin{solution}
$T_2/T_1 = (4r/r)^{3/2} = 4^{3/2} = 8$.
\end{solution}
\end{question}

\begin{question}{25} \author{J. Halliday}
\shout{Check numeric answer against the solutions key before final print.}
A rigid rod of length $L$ and mass $m$ is pivoted at one end and released
from rest in a horizontal position.

\part Write down the moment of inertia of the rod about the pivot.
\partmarks*{5}
\[
I = \tfrac{1}{3}mL^2.
\]

\part Using energy conservation, find the angular speed $\omega$ of the
rod as it passes through the vertical position.
\partmarks{12}
\begin{solution}
The centre of mass falls a height $L/2$, so $mg(L/2) = \tfrac12 I \omega^2$,
giving $\omega = \sqrt{3g/L}$.
\end{solution}

\part Find the linear speed of the free end of the rod at that instant.
\partmarks{8}
\begin{solution}
$v = \omega L = \sqrt{3gL}$.
\end{solution}
\end{question}

\end{document}
